% ============================================================
% Sección 10: Sistemas multiagente
% ============================================================
\section{Sistemas multiagente}

\begin{frame}{¿Por qué usar varios agentes?}
  \begin{itemize}
    \item Un solo agente con demasiadas responsabilidades tiende a
      \textbf{perder foco} (contexto saturado, instrucciones contradictorias).
    \item Dividir el trabajo entre agentes \textbf{especializados}
      —cada uno con su propio prompt, herramientas y contexto— suele dar
      mejores resultados en tareas grandes.
    \item Costo: más complejidad, más latencia, más tokens.
  \end{itemize}
\end{frame}

\begin{frame}{Topología 1: orquestador–trabajadores}
  \centering
  \begin{tikzpicture}[node distance=1.6cm]
    \node[cajaAccento] (orq) {Orquestador};
    \node[cajaAgente, below left=1.4cm and 1.6cm of orq] (w1) {Trabajador 1};
    \node[cajaAgente, below=1.6cm of orq] (w2) {Trabajador 2};
    \node[cajaAgente, below right=1.4cm and 1.6cm of orq] (w3) {Trabajador 3};
    \draw[flecha] (orq) -- (w1);
    \draw[flecha] (orq) -- (w2);
    \draw[flecha] (orq) -- (w3);
    \draw[flechaPunteada] (w1) to[bend right=15] (orq);
    \draw[flechaPunteada] (w3) to[bend left=15] (orq);
  \end{tikzpicture}
  \vfill
  \footnotesize Un agente central descompone la tarea, la reparte y
  \textbf{combina los resultados}. Ideal cuando las subtareas son
  independientes entre sí.
\end{frame}

\begin{frame}{Topología 2: pipeline secuencial}
  \centering
  \begin{tikzpicture}[node distance=1.8cm]
    \node[cajaAgente] (a1) {Agente\\investigador};
    \node[cajaAgente, right=of a1] (a2) {Agente\\redactor};
    \node[cajaAgente, right=of a2] (a3) {Agente\\revisor};
    \draw[flecha] (a1) -- (a2);
    \draw[flecha] (a2) -- (a3);
  \end{tikzpicture}
  \vfill
  \footnotesize Cada agente recibe la salida del anterior como entrada.
  Bueno para procesos con \textbf{etapas claras} (investigar $\to$ escribir
  $\to$ corregir).
\end{frame}

\begin{frame}{Topología 3: colaborativa / debate}
  \centering
  \begin{tikzpicture}[node distance=2.4cm]
    \node[cajaAgente] (a1) {Agente A};
    \node[cajaAgente, right=of a1] (a2) {Agente B};
    \node[cajaAccento, below=1.2cm of $(a1)!0.5!(a2)$] (j) {Juez / síntesis};
    \draw[flechaDoble] (a1) -- (a2);
    \draw[flecha] (a1) -- (j);
    \draw[flecha] (a2) -- (j);
  \end{tikzpicture}
  \vfill
  \footnotesize Los agentes discuten/critican propuestas entre sí antes de
  llegar a una respuesta final — útil para \textbf{decisiones que se
  benefician de puntos de vista distintos}.
\end{frame}

\begin{frame}{Riesgos propios de multiagente}
  \begin{itemize}
    \item \textbf{Costos y latencia} se multiplican por cada agente involucrado.
    \item \textbf{Errores en cascada}: un agente con información incorrecta
      contamina a los que dependen de él.
    \item Coordinación: decidir \emph{cuándo termina} cada agente y qué
      formato usan para comunicarse entre sí.
    \item Regla práctica: usar multiagente solo cuando un solo agente
      \textbf{ya demostró} no ser suficiente.
  \end{itemize}
\end{frame}
